\ifdefined\gkver\else\def\gkver{student}\fi
\documentclass[\gkver]{gaokaozhenti}
\begin{document}

\chapter{1999年广东}

\begin{enumerate}
\item
%% number: 1
%% paperName: 1999年广东高考第 1 题 4 分
%% typeId: 2
%% type: 多选题
%% chapter: 光学
%% point: 光的干涉
%% method: 无
%% score: 4
%% degree: 850
%% duplicateId: 0
%% body:
在下列现象中说明光具有波动性的是\xzanswer{BC}
\fourchoices[answer=BC]
{光的直线传播}
{光的衍射}
{光的干涉}
{光电效应}
%% answer: BC
%% memo:
\memoanswer{
本题原卷及材料中未提供详解，待补充。
}

\item
%% number: 2
%% paperName: 1999年广东高考第 2 题 4 分
%% typeId: 2
%% type: 多选题
%% chapter: 运动和力
%% point: 牛顿第三定律
%% method: 无
%% score: 4
%% degree: 750
%% duplicateId: 0
%% body:
汽车拉着拖车水平道路上沿直线加速行驶，根据牛顿运动定律可知\xzanswer{BC}
\fourchoices[answer=BC]
{汽车拉拖车的力大于拖车拉汽车的力}
{汽车拉拖车的力等于拖车拉汽车的力}
{汽车拉拖车的力大于拖车受到的阻力}
{汽车拉拖车的力等于拖车受到的阻力}
%% answer: BC
%% memo:
\memoanswer{
本题原卷及材料中未提供详解，待补充。
}

\item
%% number: 3
%% paperName: 1999年广东高考第 3 题 4 分
%% typeId: 2
%% type: 多选题
%% chapter: 原子和原子核
%% point: 半衰期
%% method: 无
%% score: 4
%% degree: 750
%% duplicateId: 0
%% body:
下列说法正确的是\xzanswer{BCD}
\fourchoices[answer=BCD]
{“原子由电子带正电的物质组成”是通过卢瑟福$\alpha$粒子散射实验判定的}
{玻尔理论认为原子只能处在能量不连续的一系列状态}
{放射性元素的半衰期与温度、压强无关}
{同一元素的两种同位素，其原子核内的质子数相同而中子数不同}
%% answer: BCD
%% memo:
\memoanswer{
本题原卷及材料中未提供详解，待补充。
}

\item
%% number: 4
%% paperName: 1999年广东高考第 4 题 4 分
%% typeId: 2
%% type: 多选题
%% chapter: 原子和原子核
%% point: 人工核反应
%% method: 无
%% score: 4
%% degree: 950
%% duplicateId: 0
%% body:
已知核反应的反应物，平衡反应方程时应遵循的原则有\xzanswer{CD}
\fourchoices[answer=CD]
{质子数守恒}
{中子数守恒}
{电荷数守恒}
{质量数守恒}
%% answer: CD
%% memo:
\memoanswer{
本题原卷及材料中未提供详解，待补充。
}

\item
%% number: 5
%% paperName: 1999年广东高考第 5 题 4 分
%% typeId: 1
%% type: 单选题
%% chapter: 运动和力
%% point: 动量定理
%% method: 无
%% score: 4
%% degree: 650
%% duplicateId: 0
%% body:
物体在恒力$F$作用下作直线运动，在时间$\Delta t_{1}$内速度由 0 增大到$v$，在时间$\Delta t_{2}$内速度增大到$2v$。设$F$在$\Delta t_{1}$做的功是$W_{1}$，冲量是$I_{1}$，在$\Delta t_{2}$做的功是$W_{2}$，冲量是$I_{2}$。那么\xzanswer{D}
\fourchoices[answer=D]
{$I_{1}<I_{2}$，$W_{1}=W_{2}$}
{$I_{1}<I_{2}$，$W_{1}<W_{2}$}
{$I_{1}=I_{2}$，$W_{1}=W_{2}$}
{$I_{1}=I_{2}$，$W_{1}<W_{2}$}
%% answer: D
%% memo:
\memoanswer{
本题原卷及材料中未提供详解，待补充。
}

\item
%% number: 6
%% paperName: 1999年广东高考第 6 题 4 分
%% typeId: 1
%% type: 单选题
%% chapter: 电路
%% point: 串并联电路
%% method: 无
%% score: 4
%% degree: 550
%% duplicateId: 0
%% body:
一盒子上有两个接线柱 a 和 b，与盒内电路相连。把电压表（内阻很大）接在两接线柱之间，其读数是$3\UV$；把电流表（内阻可忽略）接在两接线柱之间，其读数是$3\UA$。则下列电路中（电源内阻不计）可能为盒内电路的是\xzanswer{A}
\fourchoices[answer=A, ispicture=true, h=2.4cm]
{figs/06a.png}
{figs/06b.png}
{figs/06c.png}
{figs/06d.png}
%% answer: A
%% memo:
\memoanswer{
本题原卷及材料中未提供详解，待补充。
}

\item
%% number: 7
%% paperName: 1999年广东高考第 7 题 4 分
%% typeId: 2
%% type: 多选题
%% chapter: 曲线运动
%% point: 人造地球卫星
%% method: 无
%% score: 4
%% degree: 450
%% duplicateId: 0
%% body:
在圆轨道上运动的质量为$m$的人造地球卫星，它到地面的距离等于地球半径$R$。地面上的重力加速度为$g$，则\xzanswer{BD}
\fourchoices[answer=BD]
{卫星运动的速度为$\sqrt{2Rg}$}
{卫星运动的周期为$4\pi\sqrt{\frac{2R}{g}}$}
{卫星运动的加速度为$\frac{1}{4}g$}
{卫星的动能为$\frac{1}{4}mgR$}
%% answer: BD
%% memo:
\memoanswer{
本题原卷及材料中未提供详解，待补充。
}

\item
%% number: 8
%% paperName: 1999年广东高考第 8 题 4 分
%% typeId: 1
%% type: 单选题
%% chapter: 电磁感应
%% point: 电磁感应电路分析
%% method: 无
%% score: 4
%% degree: 450
%% duplicateId: 0
%% body:
如图所示，MN、PQ 为两平行导轨，M、P 中有一阻值为$R$的电阻，导轨处于匀强磁场中，磁感强度为$B$，磁场方向与导轨所在平面垂直，图中磁场垂直于纸面向里。有一金属圆环沿两导轨滑动，速度为$v$，与导轨接触良好，圆环的直径$d$与两导轨间的距离相等。设金属圆环的电阻可忽略。当金属环向右作匀速运动时，\xzanswer{B}
\onepicture[w=6cm]{figs/08.png}
\fourchoices[answer=B]
{有感应电流通过电阻$R$，大小为$\frac{\pi Bdv}{R}$}
{有感应电流通过电阻$R$，大小为$\frac{Bdv}{R}$}
{有感应电流通过电阻$R$，大小为$\frac{2Bdv}{R}$}
{无感应电流通过电阻$R$}
%% answer: B
%% memo:
\memoanswer{
本题原卷及材料中未提供详解，待补充。
}

\item
%% number: 9
%% paperName: 1999年广东高考第 9 题 4 分
%% typeId: 2
%% type: 多选题
%% chapter: 交流电
%% point: 变压器
%% method: 无
%% score: 4
%% degree: 650
%% duplicateId: 0
%% body:
如图所示，一理想变压器的原线圈匝数$n_{1}=1000$匝，副线圈匝数$n_{2}=200$匝，交流电源的电动势$\varepsilon=311\sin(100\pi t)\UV$，电阻$R=88\UO$。电流表和电压表对电路的影响可忽略不计。\xzanswer{AD}
\onepicture[w=8cm]{figs/09.png}
\fourchoices[answer=AD]
{A$_{1}$的示数为$0.10\UA$}
{V$_{1}$的示数为$311\UV$}
{A$_{2}$的示数为$0.75\UA$}
{V$_{2}$的示数为$44\UV$}
%% answer: AD
%% memo:
\memoanswer{
本题原卷及材料中未提供详解，待补充。
}

\item
%% number: 10
%% paperName: 1999年广东高考第 10 题 4 分
%% typeId: 1
%% type: 单选题
%% chapter: 电场
%% point: 电场叠加
%% method: 无
%% score: 4
%% degree: 550
%% duplicateId: 0
%% body:
两个固定的异号点电荷，电量给定但大小不等，用$E_{1}$和$E_{2}$分别表示两个点电荷产生的电场强度的大小，则在通过两点电荷的直线上，$E_{1}=E_{2}$的点\xzanswer{C}
\fourchoices[answer=C]
{有三个，其中两处合场强为零}
{有三个，其中一处合场强为零}
{只有两个，其中一处合场强为零}
{只有一个，该处合场强不为零}
%% answer: C
%% memo:
\memoanswer{
设异号点电荷为$-Q_{1}$和$Q_{2}$，且$Q_{1}>Q_{2}$，在通过两点电荷的直线上，两点电荷之间应有一个$E_{1}=E_{2}$处。该处的$E_{1}$、$E_{2}$方向相同。而在$Q_{2}$的外侧。也有一个$E_{1}=E_{2}$处，该处的$E_{1}$、$E_{2}$方向相反。
}

\item
%% number: 11
%% paperName: 1999年广东高考第 11 题 4 分
%% typeId: 2
%% type: 多选题
%% chapter: 力物体的平衡
%% point: 弹力
%% method: 无
%% score: 4
%% degree: 450
%% duplicateId: 0
%% body:
如图所示，a、b 为两根相连的轻质弹簧，它们的劲度系数分别为$k_{a}=1\times10^{3}\UNm$，$k_{b}=2\times10^{3}\UNm$，原长分别为$l_{a}=6\Ucm$，$l_{b}=4\Ucm$。\xzanswer{BC}
\onepicture[h=4.5cm]{figs/11.png}
\fourchoices[answer=BC]
{弹簧 a 的下端的拉力为$4\UN$，b 端受的拉力为$6\UN$}
{弹簧 a 的下端的拉力为$10\UN$，b 端受的拉力为$10\UN$}
{弹簧 a 的长度变为$7\Ucm$，b 的长度变为$4.5\Ucm$}
{弹簧 a 的长度变为$6.4\Ucm$，b 的长度变为$4.3\Ucm$}
%% answer: BC
%% memo:
\memoanswer{
本题原卷及材料中未提供详解，待补充。
}

\item
%% number: 12
%% paperName: 1999年广东高考第 12 题 4 分
%% typeId: 2
%% type: 多选题
%% chapter: 机械波
%% point: 波的图象
%% method: 无
%% score: 4
%% degree: 450
%% duplicateId: 0
%% body:
一列简谐横波沿绳子传播，振幅为$0.2\Ucm$，传播速度为$1\Ums$，频率为$0.5\UHz$。在$t_{0}$时刻，质点 a 正好经过平衡位置。沿着波的传播方向\xzanswer{BD}
\fourchoices[answer=BD]
{在$t_{0}$时刻，距 a 点$2\Um$的质点离开其平衡位置的距离为$0.2\Um$}
{在（$t_{0}+1\Us$）时刻，距 a 点为$1.5\Um$处的质点离开其平衡位置的距离为$0.2\Um$}
{在（$t_{0}+2\Us$）时刻，距 a 点为$1\Um$处的质点离开其平衡位置的距离为$0.2\Um$}
{在（$t_{0}+3\Us$）时刻，距 a 点为$0.5\Um$处的质点离开其平衡位置的距离为$0.2\Um$}
%% answer: BD
%% memo:
\memoanswer{
本题原卷及材料中未提供详解，待补充。
}

\item
%% number: 13
%% paperName: 1999年广东高考第 13 题 5 分
%% typeId: 3
%% type: 填空题
%% chapter: 磁场
%% point: 安培力
%% method: 无
%% score: 5
%% degree: 550
%% duplicateId: 0
%% body:
一劲度系数为$k$的轻质弹簧，下端挂有一匝数为$n$的矩形线圈 abcd。bc 边长为$l$。线框的下半部处在匀强磁场中，磁感强度大小为$B$，方向与线框平面垂直，在图中的垂直纸面向里。线框中通以电流$I$，方向如图所示。开始时线框处于平衡状态。令磁场反向，磁感强度的大小仍为$B$，线框达到新的平衡。在此过程中线框位移的大小$\Delta x$\tkanswer{$\frac{2nBIl}{k}$}，方向\tkanswer{向下}。
\onepicture[h=5cm, align=right]{figs/13.png}
%% answer: 
\jdanswer{
$\frac{2nBIl}{k}$，向下
}
%% memo:
\memoanswer{
本题原卷及材料中未提供详解，待补充。
}

\item
%% number: 14
%% paperName: 1999年广东高考第 14 题 5 分
%% typeId: 1
%% type: 单选题
%% chapter: 力物体的平衡
%% point: 力
%% method: 无
%% score: 5
%% degree: 950
%% duplicateId: 0
%% body:
一轻绳上端固定，下端连有一质量为、带电量为$q$的小球。现加一水平方向的匀强电场，使小球平衡时轻绳与竖直方向成$\theta$角，则此时电场强度的大小为\tkanswer{$\frac{mg\tan\theta}{q}$}。
%% answer: 
\jdanswer{
$\frac{mg\tan\theta}{q}$
}
%% memo:
\memoanswer{
本题原卷及材料中未提供详解，待补充。
}

\item
%% number: 15
%% paperName: 1999年广东高考第 15 题 5 分
%% typeId: 3
%% type: 填空题
%% chapter: 交流电
%% point: 交流电
%% method: 无
%% score: 5
%% degree: 550
%% duplicateId: 0
%% body:
一阻值恒定的电阻器，当两端加上$10\UV$的直流电压时，测得它的功率为$P$；当两端加上某一正弦交流电压时，测得它的功率为$\frac{P}{2}$。由此可知该交流电压的有效值为\tkanswer[1.2cm]{$5\sqrt{2}$}$\UV$，峰值为\tkanswer[1.2cm]{$10$}$\UV$。
%% answer: 
\jdanswer{
$5\sqrt{2}$，$10$
}
%% memo:
\memoanswer{
本题原卷及材料中未提供详解，待补充。
}

\item
%% number: 16
%% paperName: 1999年广东高考第 16 题 5 分
%% typeId: 3
%% type: 填空题
%% chapter: 曲线运动
%% point: 平抛运动
%% method: 无
%% score: 5
%% degree: 450
%% duplicateId: 0
%% body:
一导弹离地面高度为$h$水平飞行，某一时刻，导弹的速度为$v$，突然爆炸成质量相同的 A、B 两块，A、B 同时落到地面，两落地点相距$4v\sqrt{\frac{2h}{g}}$，两落地点与爆炸前导弹速度在同一竖直平面内。不计空气阻力。已知爆炸后瞬间，A 的动能$E_{kA}$大于 B 的动能$E_{kB}$，则$E_{kA}:E_{kB}=$\tkanswer[1.2cm]{$9:1$}。
%% answer: 
\jdanswer{
$9:1$
}
%% memo:
\memoanswer{
本题原卷及材料中未提供详解，待补充。
}

\item
%% number: 17
%% paperName: 1999年广东高考第 17 题 5 分
%% typeId: 4
%% type: 实验题
%% chapter: 直线运动
%% point: 游标卡尺和螺旋测微仪
%% method: 无
%% score: 5
%% degree: 950
%% duplicateId: 0
%% body:
用螺旋测微器测量一个小球的直径，结果如图所示，则球的直径是\tkanswer[1.8cm]{$10.965$}$\Umm$。
\onepicture[h=3.5cm, align=right]{figs/17.png}
%% answer: 
\jdanswer{
$10.965\Umm$
}
%% memo:
\memoanswer{
本题原卷及材料中未提供详解，待补充。
}

\item
%% number: 18
%% paperName: 1999年广东高考第 18 题 6 分
%% typeId: 4
%% type: 实验题
%% chapter: 电路
%% point: 伏安法测电阻
%% method: 无
%% score: 6
%% degree: 650
%% duplicateId: 0
%% body:
用如图所示中所给的实验器材测量一只“$12\UV$，$5\UW$”的小灯泡在不同电压下的功率，其中电流表有$3\UA$、$0.6\UA$两档，内阻可忽略，电压表有$15\UV$、$3\UV$两档，内阻很大。测量时要求加在灯泡两端的电压可连续地从$0\UV$调到$12\UV$。
\begin{enumerate}
	\item
	按要求在实物图上连线（其中部分线路已经连好）。
	\item
	某次测量时电流表的指针位置如图所示，其读数为\tkanswer[1.8cm]{$0.36$}$\UA$。
\end{enumerate}
\drawpicanswer{\onepicture[w=9cm, align=right]{figs/18a.png}}{\onepicture[w=9cm, align=right]{figs/18_an.jpg}}
\onepicture[w=4.5cm, align=right]{figs/18b.png}
%% answer: 
\jdanswer{
\begin{enumerate}
	\item
	如图所示（红线为连线方式）。
	\item
	$0.36$（写成$0.360$也给满分）
\end{enumerate}
}
%% memo:
\memoanswer{
（1）对于“$12\UV$，$5\UW$”的小灯泡其额定电流大约是$I=\frac{5}{12}\UA<0.6\UA$，安培表的量程选$0\sim0.6\UA$。根据测量要求，电压连续地从$0\UV$调到$12\UV$，应接成分压器电路，而不应接限流器电路。又因为电流内阻可忽略，电压表内阻很大对电路无影响，电流表内接或外接都可以。
}

\item
%% number: 19
%% paperName: 1999年广东高考第 19 题 6 分
%% typeId: 4
%% type: 实验题
%% chapter: 气体性质
%% point: 验证玻意耳定律
%% method: 无
%% score: 6
%% degree: 550
%% duplicateId: 0
%% body:
用“验证玻-马定律实验”的装置来测量大气压。所用注射器的最大容积为$V_{m}$，刻度全长为$L$，活塞与钩码支架的总质量为$M$，注射器被固定在竖直方向上，如图所示。在活塞两侧各悬挂 1 个质量为$m$的钩码时注射器内空气体积为$V_{1}$；除去钩码后，用弹簧秤向上拉活塞，达到平衡时注射器内空气体积为$V_{2}$，弹簧秤的读数为$F$（整个过程中，温度保持不变）。由这些数据可以求出大气压强$p_{0}=$\tkanswer{$\frac{L}{V_{m}(V_{2}-V_{1})}(MgV_{1}-MgV_{2}+2mgV_{1}+FV_{2})$}。
\onepicture[h=5cm, align=right]{figs/19.png}
%% answer: 
\jdanswer{
$\frac{L}{V_{m}(V_{2}-V_{1})}(MgV_{1}-MgV_{2}+2mgV_{1}+FV_{2})$
}
%% memo:
\memoanswer{
本题原卷及材料中未提供详解，待补充。
}

\item
%% number: 20
%% paperName: 1999年广东高考第 20 题 10 分
%% typeId: 6
%% type: 计算题
%% chapter: 光学
%% point: 透镜
%% method: 无
%% score: 10
%% degree: 650
%% duplicateId: 0
%% body:
将一薄凸透镜嵌在一遮光板中央的圆孔上，透镜边缘与圆孔边缘重合。用平行光沿主轴照射整个透镜，在透镜另一侧与透镜平行的光屏上形成一圆形光斑。若光屏与透镜间的距离为$D$，透镜的半径为$R_{1}$，光斑的半径为$R_{2}$，$R_{1}\ll R_{2}$。求此透镜的焦距$f$。
\onepicture[w=7cm, align=right]{figs/20.png}
%% answer: 
\jdanswer{
$f=\frac{R_{1}}{R_{1}+R_{2}}D$
}
%% memo:
\memoanswer{
由于凸透镜对光线有汇聚作用，所以射到图中屏上 P 点的光线是从透镜上 H 点射来的。HP 交主轴于 F 点，F 即焦点，OF 就是焦距$f$。

由 $\triangle HFO \sim \triangle PFO'$ 得
\[ \frac{f}{R_{1}}=\frac{R_{1}}{R_{1}+R_{2}} \quad \text{①} \]
\onepicture[w=7cm]{figs/20_an.jpg}

解得：$f=\frac{R_{1}}{R_{1}+R_{2}}D \quad \text{②}$
}

\item
%% number: 21
%% paperName: 1999年广东高考第 21 题 10 分
%% typeId: 6
%% type: 计算题
%% chapter: 功和能
%% point: 动能定理
%% method: 无
%% score: 10
%% degree: 550
%% duplicateId: 0
%% body:
\begin{enumerate}
	\item
	试在下列简化的情况下从牛顿运动定律出发，导出动能定理的表达式：物体为质点，作用力是恒力，运动轨道为直线。要求写出每个符号以及所得结果中每项的意义。
	\item
	如图所示，一弹簧振子，物块的质量为$m$，它与水平桌面间的动摩擦因数为$\mu$。起初，用手按住物块，物块的速度为零，弹簧的伸长量为$x$。然后放手，当弹簧的长度回到原长时，物块的速度为$v$。试用动能定理求此过程中弹力所作的功。
\end{enumerate}
\onepicture[w=7cm, align=right]{figs/21.png}
%% answer: 
\jdanswer{
\begin{enumerate}
	\item
	$W=\frac{1}{2}mv_{2}^{2}-\frac{1}{2}mv_{1}^{2}$
	\item
	$W_{\text{弹}}=\mu mgx+\frac{1}{2}mv^{2}$
\end{enumerate}
}
%% memo:
\memoanswer{
用$m$表示物体质量，$F$表示作用于物体的恒定合力，在此力作用下，物体由位置$a$沿直线移动到$b$点，设位移为$s$，物体在$a$点的速度为$v_{1}$，在$b$点速度为$v_{2}$，则外力作的功
\[ W=Fs \quad (1) \]
由牛顿第二定律
\[ F=ma \quad (2) \]
由运动学公式
\[ v_{2}^{2}=v_{1}^{2}+2as \quad (3) \]
把（2）、（3）代入（1），得到动能定理的表达式
\[ W=\frac{1}{2}mv_{2}^{2}-\frac{1}{2}mv_{1}^{2} \quad (4) \]
式中$W$为作用于物体合力所作的功，$\frac{1}{2}mv_{2}^{2}$为物体的末动能，$\frac{1}{2}mv_{1}^{2}$为物体的初动能，（4）式表示作用于物体的合力作的功等于物体动能的变化量。

（2）设$W_{\text{弹}}$为弹力对物体作的功。因为克服摩擦力作的功为$\mu mgx$，由动能定理，得：
\[ W_{\text{弹}}=-\mu mgx=\frac{1}{2}mv^{2}-0 \]
得：$W_{\text{弹}}=\mu mgx+\frac{1}{2}mv^{2}$
}

\item
%% number: 22
%% paperName: 1999年广东高考第 22 题 13 分
%% typeId: 6
%% type: 计算题
%% chapter: 功和能
%% point: 整体机械能守恒
%% method: 无
%% score: 13
%% degree: 450
%% duplicateId: 0
%% body:
如图所示，一固定的楔形木块，其斜面倾角为$\theta=30^{\circ}$，另一边与地面垂直，顶上有一个定滑轮。一柔软的细线跨过定滑轮，两端分别与物块 A 和 B 连结，A 的质量为$4m$，B 的质量为$m$。开始时将 B 按在地面上不动，然后放开手，让 A 沿斜面下滑而 B 上升。物块 A 与斜面间无摩擦。设当 A 沿斜面下滑$s$距离后，细线突然断了。求物块 B 上升的最大高度$H$。
\onepicture[w=7cm, align=right]{figs/22.png}
%% answer: 
\jdanswer{
$H=1.5s$
}
%% memo:
\memoanswer{
设物块 A 沿斜面下滑$s$距离时的速度为$v$，由机械能守恒定律得：
\[ \frac{5}{2}mv^{2}+mgs=4mgs\sin30^{\circ}=2mgs \quad (1) \]
细线突然断的瞬间，物块 B 垂直上升的速度为$v$，此后 B 作竖直上抛运动。设继续上升的距离为$h$，由机械能守恒定律得：
\[ \frac{1}{2}mv^{2}=mgh \quad (2) \]
物块 B 上升的最大高度
\[ H=h+s \quad (3) \]
由（1）、（2）、（3）得：
\[ H=1.5s \]
}

\item
%% number: 23
%% paperName: 1999年广东高考第 23 题 13 分
%% typeId: 6
%% type: 计算题
%% chapter: 气体性质
%% point: 玻意耳定律
%% method: 无
%% score: 13
%% degree: 250
%% duplicateId: 0
%% body:
麦克劳真空计是一种测量极稀薄气体压强的仪器，其基本部分是一个玻璃连通器，其上端玻璃管 A 与盛有待测气体的容器连接，其下端 D 经过橡皮软管与水银容器 R 相通，如图所示。图中 K$_{1}$、K$_{2}$是相互平行的竖直毛细管，它们的内径皆为$d$，K$_{1}$顶端封闭，在玻璃泡 B 与管 C 相通处刻有标记 m，测量时先降低 R 使水银面低于 m，如图（a）所示，逐渐提升 R，直至 K$_{2}$中水银面与 K$_{1}$顶端等高，这时 K$_{1}$中水银面比顶端低$h$，如图（b）所示。设待测容器较大，水银面升降不影响其中压强，测量过程中温度不变。已知 B（m 以上）的容积为$V$，K$_{1}$的容积远小于$V$，水银的密度为$\rho$。
\begin{enumerate}
	\item
	试导出上述过程中计算待测压强$p$的表达式；
	\item
	已知$V=628\Ucmc$，毛细管的直径$d=0.3\Umm$，水银密度$\rho=13.6\times10^{3}\Ukgmc$，$h=40\Umm$，算出待测压强$p$（计算时$g$取$10\Umsq$，结果保留两位有效数字）。
\end{enumerate}
\twopicture[labelprefix={}, numA={（a）}, numB={（b）}, h=6cm, align=right]
{figs/23a.png}
{figs/23b.png}
%% answer: 
\jdanswer{
\begin{enumerate}
	\item
	$p=\frac{\rho g h^{2}\pi d^{2}}{4V}$
	\item
	$2.4\times10^{-2}\UPa$
\end{enumerate}
}
%% memo:
\memoanswer{
（1）水银面升到 m 时，B 中气体刚被封闭，压强为待测压强$p$。这部分气体末态体积为$ah$，$a=\frac{\pi d^{2}}{4}$，压强为$p+h\rho g$，由玻意耳定律，得：
\[ pV=(p+h\rho g)\frac{\pi d^{2}}{4}h \quad \text{①} \]
整理得：$p(V-\frac{\pi d^{2}}{4}h)=h\rho g\frac{\pi d^{2}}{4}h$

根据题给条件，$\frac{\pi d^{2}}{4}h$远小于$V$，得：
\[ pV=(h\rho g)\frac{\pi d^{2}}{4}h \quad \text{②} \]
化简得：$p=\frac{\rho g h^{2}\pi d^{2}}{4V} \quad \text{③}$
\onepicture[h=5.5cm]{figs/23_an1.jpg}
\onepicture[h=5cm]{figs/23_an2.jpg}

代入数据得：$p=2.4\times10^{-2}\UPa$。
}

\item
%% number: 24
%% paperName: 1999年广东高考第 24 题 16 分
%% typeId: 6
%% type: 计算题
%% chapter: 磁场
%% point: 带电粒子在磁场中的运动
%% method: 分情况讨论
%% score: 16
%% degree: 350
%% duplicateId: 0
%% body:
如图所示，在某装置中有一匀强磁场，磁感强度为$B$，方向垂直于$Oxy$所在纸面向外。某时刻在$x=l_{0}$，$y=0$处，一质子沿$y$轴的负方向进入磁场；同一时刻，在$x=-l_{0}$，$y=0$处，一个$\alpha$粒子进入磁场，速度方向与磁场垂直。不考虑质子与$\alpha$粒子的相互作用。设质子的质量为$m$，电量为$e$。
\begin{enumerate}
	\item
	如果质子经过坐标原点 O，它的速度为多大？
	\item
	如果$\alpha$粒子与质子在坐标原点相遇，$\alpha$粒子的速度应为何值？方向如何？
\end{enumerate}
\onepicture[w=6cm, align=right]{figs/24.png}
%% answer: 
\jdanswer{
\begin{enumerate}
	\item
	$v_{p}=\frac{eBl_{0}}{2m}$
	\item
	$v_{\alpha}=\frac{\sqrt{2}eBl_{0}}{4m}$，方向有两个，用$\alpha$粒子速度方向与$x$轴正方向夹角$\theta$表示：$\theta_{1}=\frac{\pi}{4}$，$\theta_{2}=\frac{3\pi}{4}$。
\end{enumerate}
}
%% memo:
\memoanswer{
（1）根据质子进入磁场的位置和进入磁场时速度的方向，可知其圆周轨道的圆心必在$x$轴上，又因质子经过原点 O，故其轨道半径$r_{P}=\frac{1}{2}l_{0}$。设质子的速度为$v_{P}$，由牛顿定律得，
\[ \frac{mv_{P}^{2}}{r_{P}}=eBv_{P} \quad \text{①} \]
解得：$v_{P}=\frac{eBl_{0}}{2m} \quad \text{②}$

（2）质子做圆周运动的周期为
\[ T_{P}=\frac{2\pi m}{eB} \quad \text{③} \]
由于$\alpha$粒子电荷为$q_{\alpha}=2e$，质量$m_{\alpha}=4m$，故粒子做圆周运动的周期
\[ T_{\alpha}=\frac{4\pi m}{eB} \quad \text{④} \]
质子在做圆周运动的过程中，在$t=\frac{1}{2}T_{P}$，$\frac{3}{2}T_{P}$，$\frac{5}{2}T_{P}\cdots$各时刻通过 O 点，$\alpha$粒子如与质子在 O 点相遇，必在同一时刻到达 O 点。这些时刻分别对应$\frac{1}{4}T_{\alpha}$，$\frac{3}{4}T_{\alpha}$，$\cdots$。如果$\alpha$粒子在$t=\frac{1}{4}T_{\alpha}$到达 O 点，它运行了$\frac{1}{4}$周期。如在$t=\frac{3}{4}T_{\alpha}$到达 O 点，它运行了$\frac{3}{4}$周期。由此可知，$\alpha$粒子进入磁场处与 O 点之间的连线必为$\frac{1}{4}$圆周或$\frac{3}{4}$圆周所对的弦，如图（实际上$t=\frac{5}{4}T_{\alpha}$等情形不必再考虑）。进而得出粒子的轨道半径
\[ r_{\alpha}=\frac{\sqrt{2}}{2}l_{0} \quad \text{⑤} \]
\onepicture[w=6cm]{figs/24_an.jpg}

设$\alpha$粒子的速度为$v_{\alpha}$，则由牛顿定律得
\[ \frac{m_{\alpha}v_{\alpha}^{2}}{r_{\alpha}}=q_{\alpha}Bv_{\alpha} \]
注意到$m_{\alpha}=4m$，$q_{\alpha}=2e$。由⑤式得
\[ v_{\alpha}=\frac{2eBl_{0}}{4m} \quad \text{⑥} \]
但方向可有两个，用$\alpha$粒子速度方向与$x$轴正方向夹角$\theta$表示
\[ \theta_{1}=\frac{\pi}{4} \quad \text{⑦} \]
\[ \theta_{2}=\frac{3\pi}{4} \quad \text{⑧} \]
}

\end{enumerate}

\end{document}
